\subsection{FRM Exam 2007: skewness of a lognormal distribution}

\textbf{Enunciado}

\emph{The skew of a lognormal distribution is always}
\begin{enumerate}[label=\alph*.]
\item \hl{Positive}
\item Negative
\item 0
\item 3
\end{enumerate}

\subsubsection*{Solución}

Sabemos que el skewness de una lognormal es
\[
\gamma_1=(e^{\sigma^2}+2)\sqrt{e^{\sigma^2}-1},
\]
que es estrictamente positiva para toda lognormal no degenerada (\(\sigma \neq 0\)).

Intuitivamente, si \(X\) es lognormal, \(X=e^{Y}\) con \(Y\) normal. Entonces \(X>0\) y presenta cola derecha larga (asimetría positiva). Por ello su skewness es estrictamente positiva.
